\subsection{Formatting JSON inside Java}\label{sec:FormattingJSONInsideJava}
Here are some basic rules for the formatting of JSON\index{JSON}:
\begin{itemize}
	\item{All strings must be wrapped in double quotes.}
	\item{Numbers are always printed with their decimal part, even when it’s 0.}
	\item{Booleans are printed as \verb#true# and \verb#false#.}
	\item{\lstinline|NULL| is printed as \verb#null#.}
	\item{Try to keep an array on a single line, otherwise it should be split with each item on a line of its own and indented by one level (four spaces).}
	\item{Try to keep an object on a single line, otherwise, it should be split with each field/value pair on a line of its own, again indented by one level.}
	\item{For an object, the field name and the relative value should never end up on different lines. If the value is an object/array that needs to be split, its open parentheses must be kept on the same line as the field name.}
	\item{Neither objects nor arrays should have a trailing comma since it’s not allowed in JSON.}
	\item{When objects and arrays are displayed on a single line there should be a single whitespace separating the parentheses and their items.}
\end{itemize}

\keeplisting{17}
These should be sufficient for our purpose; a sample would look like this:
\begin{lstlisting}[language={}]
{
    name: "value",
    flag: true,
    number: 42.0,
    simpleArray: [1.0, 2.0, 3.0],
    stringArray: ["one", "two", "three"],
    complexArray: [
        { s: "value", n: 2.0 f: false },
        { s: "other", n: 3.14, f: true }
    ],
    complexObject: {
        name: "value",
        flag: false,
        longText: "This is a long text with line break!\nHere we continue that text in the second line.\nAnd this is the third line!"
    }       
}
\end{lstlisting}

When we embed this into a Java source file, it would look like below:
\begin{lstlisting}
final var jsonSnippet = """
````{
````    name: "value",
````    flag: true,
````    number: 42.0,
````    simpleArray: [1.0, 2.0, 3.0],
````    stringArray: ["one", "two", "three"],
````    complexArray: [
````        { s: "value", n: 2.0 f: false },
````        { s: "other", n: 3.14, f: true }
````    ],(|\newpage|)
````    complexObject: {
````        name: "value",
````        flag: false,
````        longText: "This is a long text with line break!\nHere we continue that text in the second line.\nAnd this is the third line!"
````    }       
````}\
````""";
\end{lstlisting}

%==============================================================================
% End of file!!
%==============================================================================
